\section{Our Approach}
\label{sec:method}
\input{workflow}

Figure~\ref{fig:pipeline} shows the \sysname workflow.
Given a C source file, \sysname extracts loop nests, analyzes their parallelization properties, constructs a structured prompt, invokes an LLM to generate Triton code, and optionally refines performance using GPU profiling feedback.

\subsection{Analysis and Prompt Construction}
\sysname calls PET~\cite{verdoolaege2012polyhedral} once per kernel to extract a polyhedral representation (iteration domains, access relations, schedules). From this single invocation, it derives all parallelization constraints in a unified JSON schema:
\emph{(i)}~for each loop dimension, whether it can be safely parallelized (by checking for cross-iteration dependences via ISL relation composition);
\emph{(ii)}~WAR (write-after-read) dependences that require array cloning or sequential execution;
\emph{(iii)}~reduction patterns (scalar accumulation + array read) mapped to Triton primitives;
\emph{(iv)}~memory access patterns (stencil, triangular, cross-phase) that affect kernel structure.
When polyhedral analysis fails (e.g., non-affine accesses), \sysname falls back to LLVM's \texttt{DependenceAnalysis}.

A single, pattern-agnostic function renders the analysis into structured text: it emits whatever constraints were found without pattern-specific templates. The prompt includes the C source, analysis facts (e.g., ``\texttt{j} is safe; \texttt{t} has cross-phase deps''), and Triton coding rules. The LLM receives facts and decides the implementation strategy.

\subsection{Generation and Iterative Refinement}
\sysname invokes Claude Sonnet~4~\cite{anthropic2025claude} and validates the generated kernel against the C reference using identical inputs and FP32 tolerance checks. On failure, errors are classified (compilation, numerical, missing barriers, low parallelism) and a targeted retry prompt is issued, up to 10 attempts with a context reset at attempt~6.

\subsection{Profiling-Guided Optimization}
\label{sec:profiling}
After correctness validation, \sysname optionally profiles the kernel with NVIDIA Nsight Compute (NCU), collecting SM throughput, memory throughput, occupancy, and cache hit rates. It classifies the bottleneck (\emph{compute-bound}, \emph{memory-bound}, or \emph{latency-bound}) and feeds the metrics back to the LLM with bottleneck-specific optimization guidance. The optimized kernel is re-validated for correctness and kept only if faster. This loop runs up to 3 iterations, with NCU results cached when the implementation is unchanged.
